\documentclass{article}
\usepackage[T1]{fontenc}
\usepackage{lmodern}
\usepackage{graphicx}
\usepackage{amsmath,amssymb}
\usepackage[a4paper,margin=2cm]{geometry}
\usepackage[absolute,overlay]{textpos}
\setlength{\TPHorizModule}{1mm}
\setlength{\TPVertModule}{1mm}
\usepackage[indonesian]{babel}
\usepackage{hyperref}
\setlength{\emergencystretch}{2em}
% unit: Halaman 5

\begin{document}

% === Halaman 5 (llm) ===

\begin{itemize}
    \item Setiap blok plainteks dienkripsi dalam 16 putaran \textit{enciphering}.
    \item Setiap putaran menggunakan kunci internal (kunci putaran) berbeda.
    \item Kunci internal sepanjang 48-bit dibangkitkan dari kunci eksternal
    \item Setiap blok plainteks mengalami permutasi awal (\textit{IP}), 16 putaran \textit{enciphering}, dan inversi permutasi awal (\textit{IP}^{-1}). (lihat Gambar 1)
\end{itemize}

\end{document}
